\section{Multiple Precision Integers}\label{sec:multiple-precision-integers}

This section introduces the native Isabelle-LLVM library by Spinei with contributions by Lammich \cite{spinei_isabelle_big_integer}.
It then describes the extensions I have made to use the library for Pastèque.

\subsection{Native Arbitrary Precision Integers for Isabelle-LLVM}\label{sec:native-arbitrary-precision-integers-for-isabelle-llvm}

The native precision integer library models natural numbers by lists of 64-bit machine words.
Individual 64-bit machine words are referred to as \textit{limbs}.
An abstraction function $\iexp{big\_int\_}\alpha \dbcolon 64 \app \iexp{word} \app \iexp{list} \rightarrow \iexp{nat}$ is recursively defined as
\begin{gather*}
	\iexp{big\_int\_}\alpha \app xs =
	\begin{cases}
		0                                                                    & \text{, if $xs=[ \app ]$}            \\
		\iexp{unat} \app x + 2^{64} \cdot (\iexp{big\_int\_}\alpha \app xs') & \text{, if $xs = (x \iexp{\#} xs')$}
	\end{cases}
\end{gather*}
where $\iexp{unat} \dbcolon \tv{l} \app \iexp{word} \rightarrow \iexp{nat}$ is the natural number encoded by an (unsigned) machine word of bit-width $\tv{l}$.

The library also defines an invariant $\iexp{big\_int\_invar} \dbcolon 64 \app \iexp{word} \app \iexp{list} \rightarrow \iexp{bool}$ with $\iexp{big\_int\_invar} \app xs = (xs = [ \app ] \lor \iexp{last} \app xs \neq 0)$, where $\iexp{last} \dbcolon \tv{a} \app \iexp{list} \rightarrow \tv{a}$ yields the last element of a list.
Note that $\iexp{last} \app xs \neq 0$ implies $xs$ to be non empty.

A relation $\iexp{big\_int\_rel} \dbcolon (64 \app \iexp{word} \app \iexp{list} \times \iexp{nat}) \app \iexp{set}$ is then defined such that $(ws,n) \in \iexp{big\_int\_rel} \Leftrightarrow n = \iexp{big\_int\_}\alpha \app ws \land \iexp{big\_int\_invar} \app ws$.
Note that the invariant $\iexp{big\_int\_invar}$ forces a unique representation; the lists $[\iexp{0x1},\iexp{0x0},\iexp{0x0},\iexp{0x0}]$ and $[\iexp{0x1}]$ both evaluate to $1$ under $\iexp{big\_int\_}\alpha$, but only $[\iexp{0x1}]$ satisfies the invariant.
Also, $[ \app ]$ is the unique representation of $0$.

Relations like $\iexp{big\_int\_rel}$, which relate concrete values (in this case lists of 64-bit machine words) to abstract values (in this case natural numbers), based on an abstraction function (here $\iexp{big\_int\_}\alpha$) and an invariant constraining the allowed concrete values (here $\iexp{big\_int\_invar}$) are commonly used in the refinement framework.
I refer to such relations as \textit{refinement relations} since they have a similar purpose as the refinement assertions introduced in \refsec{sec:isabelle-llvm}.
Other than refinement assertions however, refinement relations do not require the concrete values to be representable in LLVM.
The concrete values of a refinement relation, in this case $64 \app \iexp{word} \app \iexp{list}$ can serve as an intermediate abstraction between the final $\iexp{llM}$-level representation, which in case of the multiple precision library is an array of 64-bit integers refining $64 \app \iexp{word} \app \iexp{list}$, and the abstract type, in this case $\iexp{nat}$.
Refinement assertions and refinement relations can be composed to obtain a full refinement assertion defining refinement of the high level HOL-type (in this case $\iexp{nat}$) and the LLVM-representable final refinement target (in this case arrays of unsigned 64-bit integers).
The multiple precision integer library for instance defines $\iexp{bi\_assn} \dbcolon \iexp{nat} \rightarrow 64 \app \iexp{word} \times 64 \app \iexp{word} \times 64 \app \iexp{word} \app \iexp{ptr} \rightarrow \iexp{assn}$ through a composition of an auxiliary refinement assertion $\iexp{bi\_aux\_assn} \dbcolon 64 \app \iexp{word} \app \iexp{list} \rightarrow 64 \app \iexp{word} \times 64 \app \iexp{word} \times 64 \app \iexp{word} \app \iexp{ptr} \rightarrow \iexp{assn}$ and $\iexp{big\_int\_rel}$, where $64 \app \iexp{word} \times 64 \app \iexp{word} \times 64 \app \iexp{word} \app \iexp{ptr}$ is the $\iexp{llM}$-level representation of a resizable array of 64-bit machine words (the first two entries carry the array's size and capacity).

Spinei implements addition, subtraction and multiplication operations for the 64-bit word lists, including correctness proofs of the form
\begin{gather*}
	(xs,n) \in \iexp{big\_int\_rel} \land (ys,m) \in \iexp{big\_int\_rel} \\
	\Rightarrow (\iexp{big\_int\_add} \app xs \app ys, n + m) \in \iexp{big\_int\_rel}
\end{gather*}
where $\iexp{big\_int\_add} \dbcolon 64 \app \iexp{word} \app \iexp{list} \rightarrow 64 \app \iexp{word} \app \iexp{list} \rightarrow 64 \app \iexp{word} \app \iexp{list}$ is the addition implementation.
The other operations have analogous correctness proofs.
The library contains a work-in-progress implementation of Karatsuba multiplication \cite{karatsuba} and a completed school book multiplication.
For Pastèque, I use the school book multiplication.

The sepref tool then synthesizes an $\iexp{llM}$-level refinement of the $\iexp{nres}$-level arithmetic functions like $\iexp{big\_int\_add}$.
This introduces a theoretical upper bound on the length of the arrays, as the length of arrays are stored in a (signed) 64-bit integer meaning that numbers $\geq (2^{64})^{2^{63}-1}$ can not be represented.
In addition and multiplication operations, this bound can be exceeded theoretically.
Ideally, a dynamic runtime check would abort the program at runtime, if the upper bound is exceeded.
At the time of writing this thesis, this is not supported in Isabelle-LLVM yet (it is planned to add this feature to Isabelle-LLVM in the future).
To still obtain a usable multiple precision integer implementation without a dynamic overflow check, the library assumes that operations do not overflow using Isabelle's \inlinecode{sorry}.
This is the only position in the library, and in the Isabelle-LLVM port of Pastèque, where \inlinecode{sorry} is used.

\subsection{Signed Arbitrary Precision Integers}\label{sec:signed-arbitrary-precision-integers}

Spinei's library refines $\iexp{nat}$ by lists of 64-bit machine words.
Pastèque uses (signed) multiple precision integers for polynomial coefficients, meaning an unbounded refinement target for Isabelle's $\iexp{int}$ type is required.
I therefore extend the library by representing signed integers by tuples of type $(64 \app \iexp{word} \app \iexp{list} \times \iexp{bool})$, where the list represents the number's absolute value and the boolean represents the sign.
Analogous to the library, I define an abstraction function $\iexp{signed\_big\_int\_}\alpha \dbcolon (64 \app \iexp{word} \app \iexp{list} \times \iexp{bool}) \rightarrow \iexp{int}$ as
\begin{gather*}
	\iexp{signed\_big\_int\_}\alpha (ws,\sigma) =
	\begin{cases}
		- \iexp{int} \app (\iexp{big\_int\_}\alpha \app ws) & \text{, if $\sigma$} \\
		\iexp{int} \app (\iexp{big\_int\_}\alpha \app ws)   & \text{, otherwise}
	\end{cases}
\end{gather*}
where $\iexp{int}$ converts $\iexp{nat}$ to $\iexp{int}$.
I also define an invariant $\iexp{signed\_big\_int\_invar} \dbcolon (64 \app \iexp{word} \app \iexp{list} \times \iexp{bool}) \rightarrow \iexp{bool}$ as
\begin{gather*}
	\iexp{signed\_big\_int\_invar} \app (ws,\sigma) = \iexp{big\_int\_invar} \app ws \land (ws=[ \app ] \Rightarrow \lnot \sigma)
\end{gather*}
which forces a unique representation of 0, since $\iexp{signed\_big\_int\_invar} \app ([ \app ],\iexp{False}) = \iexp{True}$, while $\iexp{signed\_big\_int\_invar} \app ([ \app ], \iexp{True}) = \iexp{False}$.
I then define a relation $\iexp{signed\_big\_int\_rel} \dbcolon ((64 \app \iexp{word} \app \iexp{list} \times \iexp{bool}) \times \iexp{int}) \app \iexp{set}$ with
\begin{gather*}
	((ws,\sigma),i) \in \iexp{signed\_big\_int\_rel} \\
	\Leftrightarrow i = \iexp{signed\_big\_int\_}\alpha (ws,\sigma) \land \iexp{signed\_big\_int\_invar} (ws,\sigma)
\end{gather*}
analogously to $\iexp{big\_int\_rel}$.

Implementing arithmetic functions for addition and multiplication, which are the only operations needed for Pastèque, is straightforward by case analysis on the signs of two given input numbers.
Addition of two inputs $(ws_1,\sigma_1), (ws_2,\sigma_2)$ is achieved by
\begin{align*}
	 & \iexp{signed\_big\_int\_add} \app (ws_1,\sigma_1) \app (ws_2,\sigma_2)                                                                                                     \\
	 & = \begin{cases}
		     (\iexp{big\_int\_add} \app ws_1 \app ws_2,\iexp{False}) & \text{, if $\lnot \sigma_1 \land \lnot \sigma_2$}                                                              \\
		     (\iexp{big\_int\_add} \app ws_1 \app ws_2,\iexp{True})  & \text{, if $\sigma_1 \land \sigma_2$}                                                                          \\
		     (\iexp{big\_int\_sub} \app ws_1 \app ws_2,\sigma_1)     & \text{, if $\iexp{big\_int\_leq} \app ws_2 \app ws_1$} \\
		     (\iexp{big\_int\_sub} \app ws_2 \app ws_1,\sigma_2)     & \text{, otherwise}
	     \end{cases}
\end{align*}
where $\iexp{big\_int\_sub}$ is the subtraction implementation from the library and $\iexp{big\_int\_leq} \app ws_2 \app ws_1$ is the implementation of $\leq$ for lists of limbs.
Multiplication is implemented by
\begin{gather*}
	\iexp{signed\_big\_int\_mult} \app (ws_1,\sigma_1) \app (ws_2,\sigma_2) = (\iexp{big\_int\_mult} \app ws_1 \app ws_2, \sigma_1 \iexp{XOR} \sigma_2)
\end{gather*}
where $\iexp{big\_int\_mult}$ is the multiplication implementation from the library.

I then prove correctness analogously to the unsigned case; in case of addition for example, I prove
\begin{gather*}
	((ws_1,\sigma_1),i_1) \in \iexp{signed\_big\_int\_rel} \land ((ws_2,\sigma_2),i_2) \in \iexp{signed\_big\_int\_rel} \\
	\Rightarrow (\iexp{signed\_big\_int\_add} \app (ws_1,\sigma_1) \app (ws_2,\sigma_2), i_1 + i_2) \in \iexp{signed\_big\_int\_rel}
\end{gather*}
For multiplication I prove an analogous lemma.

Unwrapping the tuple at $\iexp{llM}$-level would usually destroy the original tuple, similarly to the $\iexp{take\_impl}$ example in \refsec{sec:ownership-in-isabelle-llvm}.
I use contributions by Fleury to achieve non-destructive arithmetic operations through extracting the pure sign flag from the tuple.

\subsection{Relating Integers and ASCII Encoded Strings}\label{sec:relating-integers-and-ascii-encoded-strings}

For parsing input files, and for printing error messages, Pastèque requires conversion between arbitrary precision integers and ASCII encoded strings.
I implement end-to-end verified conversion functions, using schoolbook division for printing multiple precision integers.
For the correctness proof, I define a relation between Isabelle/HOL's $\iexp{string}$ and $\iexp{int}$ types to define semantics of integers as ASCII encoded strings.

The digits $"0"-"9"$ have ASCII codes $48-57$; I define conversion from Isabelle/HOL's $\iexp{char}$ type to $\iexp{nat}$ by a function
\begin{gather*}
	\iexp{char\_uval} \app c = (\iexp{of\_char} \app c) - 48
\end{gather*}
where $\iexp{of\_char} \dbcolon \iexp{char} \rightarrow \iexp{nat}$ converts a character to its ASCII code.
I also define an invariant $\iexp{is\_ascii\_char} \dbcolon \iexp{char} \rightarrow \iexp{bool}$\footnote{In the actual theory \isaname{LLVM_ASCII_String}, the invariant is named $\iexp{is\_ascii\_unum}$.} by
\begin{gather*}
	\iexp{is\_ascii\_char} \app c = 48 \leq (\iexp{of\_char} \app c) \land (\iexp{of\_char} \app c) \leq 57.
\end{gather*}
A relation between ASCII digits and natural numbers is then given by relation $\iexp{ascii\_char\_nat\_rel} \dbcolon (\iexp{char} \times \iexp{nat}) \app \iexp{set}$ where
\begin{gather*}
	(c,n) \in \iexp{ascii\_char\_nat\_rel} \Leftrightarrow n = \iexp{char\_uval} \app c \land \iexp{is\_ascii\_char} \app c.
\end{gather*}
The relation between ASCII characters and natural numbers is extended to a relation between ASCII strings and natural numbers through a conversion function $\iexp{str\_uval} \dbcolon \iexp{string} \rightarrow \iexp{nat}$
with
\begin{gather*}
	\iexp{str\_uval} = \iexp{foldl} \app (\lambda acc \app c. \app 10 \cdot acc + \iexp{char\_uval} \app c) \app 0.
\end{gather*}
An invariant for ASCII string representing natural numbers is given by $\iexp{is\_ascii\_unum\_str} \dbcolon \iexp{string} \rightarrow \iexp{bool}$ where
\begin{gather*}
	\iexp{is\_ascii\_unum\_str} \app cs = (cs \neq [ \app ] \land \forall c \in \iexp{set} \app cs. \app \iexp{is\_ascii\_char} \app c)
\end{gather*}
The relation $\iexp{ascii\_str\_nat\_rel} \dbcolon (\iexp{string} \times \iexp{nat}) \app \iexp{set}$ is then defined with
\begin{gather*}
	(cs,n) \in \iexp{ascii\_str\_nat\_rel} \Leftrightarrow n = \iexp{str\_uval} \app cs \land \iexp{is\_ascii\_unum\_str} \app cs.
\end{gather*}
To finally obtain a relation between ASCII string and integers, I extend $\iexp{ascii\_str\_nat\_rel}$ to consider a leading $\iexp{'-'}$ character which has ASCII code $45$.
The final conversion function $\iexp{str\_sval} \dbcolon \iexp{string} \rightarrow \iexp{int}$ is implemented as
\begin{gather*}
	\iexp{str\_sval} \app cs =
	\begin{cases}
		- \iexp{int} \app (\iexp{str\_uval} \app cs') & \text{, if $cs = c \iexp{\#} cs' \land (\iexp{of\_char} \app c) = 45$} \\
		\iexp{int} \app   (\iexp{str\_uval} \app cs)  & \text{, otherwise.}
	\end{cases}
\end{gather*}
An invariant for ASCII string representations of integers is given by $\iexp{is\_ascii\_snum\_str} \dbcolon \iexp{string} \rightarrow \iexp{bool}$ with
\begin{gather*}
	\iexp{is\_ascii\_snum\_str} \app cs =
	\begin{cases}
		\iexp{is\_ascii\_unum\_str} cs' & \text{, if $cs = c \iexp{\#} cs' \land (\iexp{of\_char} \app c) = 45$} \\
		\iexp{is\_ascii\_unum\_str} cs  & \text{, otherwise.}
	\end{cases}
\end{gather*}
The final relation $\iexp{ascii\_str\_int\_rel} \dbcolon (\iexp{string} \times \iexp{int}) \app \iexp{set}$, which is used in the correctness proofs of the string to multiple precision integer conversion (parsing), is then given by
\begin{gather*}
	(cs,i) \in \iexp{ascii\_str\_int\_rel} \Leftrightarrow i = \iexp{str\_sval} cs \land \iexp{is\_ascii\_snum\_str} cs.
\end{gather*}

Conversely, I define a relation $\iexp{int\_ascii\_str\_rel} \dbcolon (\iexp{int} \times \iexp{string}) \app \iexp{set}$ for verifying the multiple precision integer to string conversion (printing).
Unlike $\iexp{ascii\_str\_int\_rel}$, $\iexp{int\_ascii\_str\_rel}$ does not require defining an invariant on the integers, as every integer must be representable by a string, while not every string can be converted to an integer.
The invariant $\iexp{is\_ascii\_snum\_str}$ is satisfied exactly by strings which can be parsed to an integer.

Conversion from natural numbers to strings is implemented by $\iexp{chars\_of\_nat} \dbcolon \iexp{nat} \rightarrow \iexp{string}$ with
\begin{gather*}
	\iexp{chars\_of\_nat} \app n =
	\begin{cases}
		[\iexp{char\_of} (n + 48)]                                                                                                    & \text{, if $n<10$} \\
		\iexp{chars\_of\_nat} (\lfloor \frac{n}{10} \rfloor) @ [\iexp{char\_of} \app (n \text{ mod } 10 + 48)] & \text{, otherwise}
	\end{cases}
\end{gather*}
where $\iexp{char\_of} \dbcolon \iexp{nat} \rightarrow \iexp{char}$ converts a natural number $n$ to the character with ASCII code $n \text{ mod } 256$.
Extending to conversion from integers to strings is achieved by $\iexp{chars\_of\_int} \dbcolon \iexp{int} \rightarrow \iexp{string}$ where
\begin{gather*}
	\iexp{chars\_of\_int} \app i =
	\begin{cases}
		\iexp{chars\_of\_nat} \app (\iexp{nat} \app (\iexp{abs} \app i))                          & \text{, if $0 \leq i$} \\
		(\iexp{'-'}) \iexp{\#} (\iexp{chars\_of\_nat} \app (\iexp{nat} \app (\iexp{abs} \app i))) & \text{, otherwise.}
	\end{cases}
\end{gather*}
The function $\iexp{abs} \dbcolon \iexp{int} \rightarrow \iexp{int}$ returns the absolute value of a given integer.
The relation $\iexp{int\_ascii\_str\_rel} \dbcolon (\iexp{int} \times \iexp{string}) \app \iexp{set}$ is then given by
\begin{gather*}
	(i,cs) \in \iexp{int\_ascii\_str\_rel} \Leftrightarrow cs = \iexp{chars\_of\_int} \app i
\end{gather*}
As a corollary, I prove that $\iexp{int\_ascii\_str\_rel} \circ \iexp{ascii\_str\_int\_rel} = \iexp{Id}$, where $\circ \dbcolon (\tv{a} \times \tv{b}) \app \iexp{set} \rightarrow (\tv{b} \times \tv{c}) \app \iexp{set} \rightarrow (\tv{a} \times \tv{c}) \app \iexp{set}$ is the composition of two relations and $\iexp{Id} \dbcolon (\tv{a} \times \tv{a}) \app \iexp{set}$ is the identity relation.
Note that the inverse does not hold, i.e. $\iexp{ascii\_str\_int\_rel} \circ \iexp{int\_ascii\_str\_rel} \neq \iexp{Id}$, because $\iexp{ascii\_str\_int\_rel}$ is not injective as the strings $\iexp{"01"}$ and $\iexp{"1"}$ are both mapped to $1$.

Note that the definition of $\iexp{chars\_of\_nat}$ indicates the approach I chose to implement conversion of integer to string: a given number is continuously divided by 10 (with remainder) until the result is representable by a single digit; the resulting string is then constructed from the collected remainders and final digit.

\subsection{Implementation of Conversion between String and Multiple Precision Integer}\label{sec:conversion-between-string-and-integer}
The previous subsection, \refsec{sec:relating-integers-and-ascii-encoded-strings}, defines the semantics of integers encoded as ASCII strings.
This subsection describes the concrete implementation of the conversion between strings and arbitrary precision integers, required for parsing input files and printing coefficients in error messages.

Pastèque uses the high level $\iexp{string} = \iexp{char} \app \iexp{list}$ type for strings, in particular for error messages.
Isabelle-LLVM does not define a refinement target for $\iexp{char}$, so a custom refinement assertion is required for this thesis.
The high level $\iexp{char}$ type is defined by eight booleans, i.e. $\iexp{char} = \iexp{Char} \app (\iexp{digit0} \dbcolon \iexp{bool}) \app \dots \app (\iexp{digit7} \dbcolon \iexp{bool})$, so the natural refinement target is $8 \app \iexp{word}$.
I use the existing refinement assertion $\iexp{un8\_assn} \dbcolon \iexp{nat} \rightarrow 8 \app \iexp{word} \rightarrow \iexp{assn}$\footnote{I simplify the name here, the actual refinement assertion is $\iexp{unat\_assn' TYPE(8)}$.}, refining natural numbers $< 256$ by 8-bit machine words and compose it with the existing $\iexp{char\_of} \dbcolon \iexp{nat} \rightarrow \iexp{char}$ function to obtain $\iexp{char\_assn} \dbcolon \iexp{char} \rightarrow 8 \app \iexp{word} \rightarrow \iexp{assn}$.
I then use the list implementation, discussed in \refsec{sec:lists-with-tail-pointers}, to obtain a refinement assertion $\iexp{strl\_assn}$ for strings.
\todo{Maybe mention the alternative stra\_assn in the finalization section?}

Implementing conversion between strings and integers amounts to refining the functions $\iexp{chars\_of\_int}$ and $\iexp{str\_sval}$.
For $\iexp{str\_sval}$ this is straightforward by use of the $\iexp{foldl}$ implementation for lists in LLVM which is discussed in \refsec{sec:non-destructive-list-access-using-fold}.
Refinement of $\iexp{chars\_of\_int}$ is more involved as it requires division of multiple precision integers by 10.
In particular, for a given unsigned multiple precision integer $n$, the implementation must compute $\lfloor \frac{n}{10} \rfloor $ and $n \text{ mod } 10$.

The division is implemented by classic schoolbook division.
The core of the division is implemented by a function $\iexp{bi\_div\_by\_w64}$, a simplified version of which\footnote{the actual definition in theory \isaname{LLVM_ASCII_String} also includes an invariant required for correctness and $\iexp{ASSERT}$ statements required for refinement to $\iexp{llM}$.} is shown in \refalgo{algo:bi-div-by-w64}, where $\iexp{big\_int} \equiv 64 \app \iexp{word} \app \iexp{list}$.
\isabellecodeextern[linerange=12-28,label=algo:bi-div-by-w64,caption={
			Core algorithm for dividing multiple precision numbers by a given 64-bit integer.
		}]{isabelle-code/BigIntString_Simplified.thy}
The helper function $\iexp{divmod2by1}$ is shown in \refalgo{algo:divmod2by1}.
\isabellecodeextern[linerange=6-9,label=algo:divmod2by1,caption={
			Helper function dividing a 128-bit integer by a 64-bit integer.
		}]{isabelle-code/BigIntString_Simplified.thy}
The reason that $\iexp{divmod2by1}$ is a separate function lies in the multiplication of $\iexp{hi}$ by $2^{64}$.
The limbs in the list are 64-bit integers which cannot express $\iexp{hi} \cdot 2^{64}$.
To still implement $\iexp{divmod2by1}$, the inputs $\iexp{hi}$ and $\iexp{lo}$ are upcast to a 128-bit integer before the division and modulo operations are performed.
Note that throughout the execution of $\iexp{bi\_div\_by\_w64}$ it holds that $\iexp{r} < \iexp{l}$ which implies $\iexp{cur} \app \iexp{div} \app \iexp{unat} \app \iexp{l} < 2^{64}$ in the call to $\iexp{divmod2by1}$:
\begin{itemize}
    \item $\iexp{r}$ is initialized with $0$.
    \item After a call to $\iexp{divmod2by1}$, $\iexp{r}$ is updated to $\iexp{cur} \app \iexp{mod} \app \iexp{unat} \app \iexp{l} < \iexp{unat} \app \iexp{l} \Rightarrow \iexp{unat} \app \iexp{r} + 1 \leq \iexp{unat} \app \iexp{l}$.
    \item Thus,
    \begin{align*}
        \left \lfloor \frac{\iexp{cur}}{\iexp{unat} \app \iexp{l}} \right \rfloor =
        &\left \lfloor \frac{(\iexp{unat} \app \iexp{r})\cdot2^{64} + \iexp{unat} \app \iexp{bii}}{\iexp{unat} \app \iexp{l}} \right \rfloor \\
        &\leq \frac{(\iexp{unat} \app \iexp{r})\cdot2^{64} + \iexp{unat} \app \iexp{bii}}{\iexp{unat} \app \iexp{l}} \\
        &<\frac{(\iexp{unat} \app \iexp{r} + 1)\cdot2^{64}}{\iexp{unat} \app \iexp{l}} \\
        &\leq 2^{64}.
    \end{align*}
\end{itemize}

This means that $\iexp{cur} \app \iexp{div} \app \iexp{unat} \app \iexp{l}$ can be downcast to a 64-bit integer without its value being changed.
As $\iexp{cur} \app \iexp{mod} \app \iexp{unat} \app \iexp{l} < 2^{64}$ is always fulfilled (since $\iexp{l} < 2^{64}$), it can also be downcast to a 64-bit integer without issues.
